% AUTO-GENERATED. Edit the Markdown sources or notes.config.json.
% Sources: 02-Economics/2026-Financial-Economics/2026-09-17-Note.md

\section{Financial Decisions and the Financial System / 金融与财务决策}

\subsection{The Scope of Finance / 金融学的研究对象}

\textbf{Definition}

\begin{quote}
\textbf{Finance} is the study of how people allocate scarce resources over time.
\end{quote}

金融决策有两个共同特征：

\begin{itemize}
\tightlist
\item
  成本与收益分布在不同时间；
\item
  未来现金流的发生时间与规模通常只具有概率意义，因此必须处理 \textbf{risk}。
\end{itemize}

\textbf{Definition}

\begin{quote}
The \textbf{financial system} is the set of markets and other institutions used for financial contracting and exchange of assets and risks.
\end{quote}

实施金融决策依赖 financial system。课程材料强调，它的最终功能是服务实体经济并满足人们的消费偏好。

金融理论包含两层工具：

\begin{enumerate}
\def\labelenumi{\arabic{enumi}.}
\tightlist
\item
  用于组织跨期资源配置思考的概念框架；
\item
  用于评价备选方案、作出决策并执行决策的定量模型。
\end{enumerate}

金融体系最终服务于实体经济。企业、政府与各类金融安排本身不是终点，其功能是帮助个人实现消费、储蓄和风险承担方面的偏好。

\subsection{Household Financial Decisions / 家庭金融决策}

家庭金融决策可以归纳为四类：

\begin{longtable}[]{@{}
  >{\raggedright\arraybackslash}p{(\linewidth - 2\tabcolsep) * \real{0.5000}}
  >{\raggedright\arraybackslash}p{(\linewidth - 2\tabcolsep) * \real{0.5000}}@{}}
\toprule\noalign{}
\begin{minipage}[b]{\linewidth}\raggedright
决策
\end{minipage} & \begin{minipage}[b]{\linewidth}\raggedright
核心问题
\end{minipage} \\
\midrule\noalign{}
\endhead
\bottomrule\noalign{}
\endlastfoot
Consumption and saving & 当前财富用于多少消费、保留多少用于未来？ \\
Investment / asset allocation & 储蓄如何在房地产、股票、债券等资产之间配置？ \\
Financing & 何时借入资金，以何种方式为消费或投资计划融资？ \\
Risk management & 面临哪些风险，如何降低不希望承担的财务不确定性？ \\
\end{longtable}

家庭资产负债表的基本恒等式为

\[
\text{Net Worth}=\text{Assets}-\text{Liabilities}.
\]

这里的 \textbf{asset allocation} 是对已有储蓄存量的配置；\textbf{financing decision} 则改变负债及可支配资金来源，两者不应混为一谈。

\subsection{Corporate Financial Decisions / 企业金融决策}

\subsubsection{Strategic Planning / 战略规划}

战略规划首先回答企业``做什么''：服务哪些客户、使用什么技术、生产哪些产品。战略会随技术进步和 \textbf{creative destruction} 调整，并为后续投资与融资决策设定边界。

\subsubsection{Capital Budgeting / 资本预算}

\textbf{Definition}

\begin{quote}
The \textbf{capital budgeting process} is the preparation of a plan for acquiring factories, machinery, research laboratories, showrooms, warehouses, and human assets to implement the strategic plan.
\end{quote}

基本分析单位是 \textbf{investment project}，过程包括项目识别、评价与实施。

\subsubsection{Financing and Capital Structure / 融资与资本结构}

\textbf{Definition}

\begin{quote}
\textbf{Capital structure} is the amount of the firm's market value allocated to each category of financial capital.
\end{quote}

项目可由留存收益、股票或债券等来源融资。Capital structure 同时影响：

\begin{itemize}
\tightlist
\item
  未来现金流在不同资本提供者之间的分配；
\item
  各方承担的风险；
\item
  不同状态下企业的控制权。
\end{itemize}

普通股股东通常选举董事会；优先股股东在股息未支付等特定情形下可能获得额外控制权；债券契约则通过 \textbf{covenants} 限制损害债权人利益的行为。

\subsubsection{Working Capital Management / 营运资本管理}

\textbf{Definition}

\begin{quote}
\textbf{Working capital} is the difference between current assets and current liabilities.
\end{quote}

\[
\text{Net Working Capital}
=\text{Current Assets}-\text{Current Liabilities}.
\]

流动资产通常在一年内转化为现金，流动负债通常在一年内到期。盈利不等于有流动性：即使长期有利可图，企业也可能因短期现金不足而陷入困境。

\subsection{Forms of Business Organization / 企业组织形式}

\begin{longtable}[]{@{}
  >{\raggedright\arraybackslash}p{(\linewidth - 6\tabcolsep) * \real{0.2500}}
  >{\raggedright\arraybackslash}p{(\linewidth - 6\tabcolsep) * \real{0.2500}}
  >{\raggedright\arraybackslash}p{(\linewidth - 6\tabcolsep) * \real{0.2500}}
  >{\raggedright\arraybackslash}p{(\linewidth - 6\tabcolsep) * \real{0.2500}}@{}}
\toprule\noalign{}
\begin{minipage}[b]{\linewidth}\raggedright
形式
\end{minipage} & \begin{minipage}[b]{\linewidth}\raggedright
法律与治理
\end{minipage} & \begin{minipage}[b]{\linewidth}\raggedright
责任
\end{minipage} & \begin{minipage}[b]{\linewidth}\raggedright
税收与延续性
\end{minipage} \\
\midrule\noalign{}
\endhead
\bottomrule\noalign{}
\endlastfoot
Sole proprietorship & 个人或家庭所有，所有者与企业不分离 & 通常为无限责任 & 管理成本低，通常按个人层面纳税 \\
Partnership & 两名及以上所有者，协议规定决策和损益分配 & General partner 通常负无限责任；limited partner 不参与日常管理 & 通常避免公司与个人两层重复征税 \\
Corporation & 独立法人，可持有财产、借款、缔约、起诉或被诉 & 股东通常承担有限责任 & 股份转让不终止企业；C corporation 可能发生双重征税 \\
\end{longtable}

课件中的餐馆例子展示了 \textbf{double taxation}：若公司先按 \(35\%\) 纳税，再由股东对股息按 \(28\%\) 纳税，则税后现金流为

\[
100{,}000(1-35\%)(1-28\%)=46{,}800,
\]

低于个人独资情形下的

\[
100{,}000(1-28\%)=72{,}000.
\]

\textbf{REIT} 是针对房地产投资的特殊制度安排。课件给出的典型要求包括股东人数、持股集中度、收入分配比例以及房地产相关资产和收入占比等；这些属于特定法域规则，实际应用时应核对最新法规。

课堂还介绍了滴灌通的 \textbf{Daily Revenue Contract, DRC}：投资回报随门店每日收入分成，而非固定利息或传统股权分红。理解这类合约时，应分别判断现金流索取权、控制权、期限、损失承担方式和监管分类，不能只凭``非股非债''的标签判断经济实质。

\subsection{Ownership, Management, and Corporate Governance / 所有权、经营权与公司治理}

\subsubsection{Why Ownership and Management Are Separated / 所有权与经营权分离的原因}

股东把经营权交给职业经理人，可能获得：

\begin{itemize}
\tightlist
\item
  专业技能与规模经济；
\item
  所有者投资组合的分散化；
\item
  信息搜集成本的节约；
\item
  组织学习与持续经营能力。
\end{itemize}

代价是 \textbf{agency conflict}：经理人的目标可能偏离所有者目标。因此，所有权与经营权分离是否有效，取决于监督、激励与问责机制能否以合理成本缓解冲突。

\subsubsection{Managerial Objectives / 管理层目标}

\textbf{Management rule}

\begin{quote}
Maximize the wealth of current shareholders.
\end{quote}

这条规则比``利润最大化''更明确，因为利润最大化至少存在两个问题：

\begin{enumerate}
\def\labelenumi{\arabic{enumi}.}
\tightlist
\item
  多期经营中，究竟最大化哪一期利润？
\item
  未来收入和成本不确定时，``利润''如何在风险不同的方案间比较？
\end{enumerate}

股价在有效市场中反映预期未来现金流的现值及其风险，因此为跨期决策提供同一尺度。但股东财富最大化并不意味着可以忽略合同、法律和其他利益相关者；员工、供应商、客户等 \textbf{stakeholders} 会影响企业长期现金流、风险与声誉。

\subsubsection{Market Discipline and Corporate Takeovers / 市场约束与公司接管}

上市公司股权分散时，单个股东通常缺少搜集信息和组织变革的激励。当经营不善压低股价，收购者可能取得控制权、更换管理层并通过价值改善获利。\textbf{Takeover market} 因而可以形成外部治理约束，但它依赖信息、融资条件、控制权市场和资本市场定价是否有效。

\subsection{Financial Functions within the Firm / 企业中的财务职能}

\begin{longtable}[]{@{}
  >{\raggedright\arraybackslash}p{(\linewidth - 2\tabcolsep) * \real{0.5000}}
  >{\raggedright\arraybackslash}p{(\linewidth - 2\tabcolsep) * \real{0.5000}}@{}}
\toprule\noalign{}
\begin{minipage}[b]{\linewidth}\raggedright
角色
\end{minipage} & \begin{minipage}[b]{\linewidth}\raggedright
主要职责
\end{minipage} \\
\midrule\noalign{}
\endhead
\bottomrule\noalign{}
\endlastfoot
CEO & 向董事会负责，统筹企业经营 \\
CFO & 对企业全部财务职能负总责 \\
Treasurer & 现金、融资、银行关系、投资者关系、汇率与利率风险、税务等 \\
VP Financial Planning & 资本支出、新业务、退出、并购与分拆分析 \\
Controller / Comptroller & 财务、管理与成本会计，审计，预算与实际比较，对外报表 \\
\end{longtable}

更完整的财务职能还包括规划、资本筹集、资金管理、会计与控制、资产保护、税务、投资者关系、评价与咨询，以及 management information systems。

\begin{quote}
\textbf{注：课堂扩展：AI 与金融职业}

CH1 pp.54--61 讨论 AI 对前中后台工作的自动化，以及 prompt engineering、AI ethics、算法合规、数字资产和人机协作等可能方向。这部分是课堂扩展，不是本章金融理论主线；其中涉及机构案例、替代率、裁员率和薪酬的数字未在课件中给出可追溯来源，引用前应另行核验。
\end{quote}
